\section*{Sources and evidence}
\addcontentsline{toc}{section}{Sources and evidence}
The equal-block construction, proof, executable search, generated tables and
illustrations are original teaching artifacts. On 25 September 2026 the
Watson--Crick automata paper's introduction, preliminaries and acceptance
definition were inspected. Our empty-input and nonempty-transition conventions
are stated explicitly; no hierarchy theorem is transferred to this variant.
The Benenson paper's mechanism description and Figure 2 caption were inspected
in the author-course archive copy. The mechanism plate is an original
conceptual abstraction, not a reproduction, sequence design or wet-lab protocol.
Computed configurations are not measured chemical trajectories.
\begingroup\let\chapter\section
\begin{thebibliography}{9}
\bibitem{czeizler} E. Czeizler, E. Czeizler, L. Kari and K. Salomaa.
\emph{On the descriptional complexity of Watson--Crick automata}.
Theoretical Computer Science 410 (2009), 3250--3260.
\url{https://doi.org/10.1016/j.tcs.2009.05.001}.
Author copy: \url{https://cs.uwaterloo.ca/~lila/pdfs/Watson_Crick_automata.pdf}.
\bibitem{benenson} Y. Benenson, T. Paz-Elizur, R. Adar, E. Keinan,
Z. Livneh and E. Shapiro.
\emph{Programmable and autonomous computing machine made of biomolecules}.
Nature 414 (2001), 430--434.
\url{https://doi.org/10.1038/35106533}.
Mechanism inspected in the Duke molecular-computing course archive copy.
\end{thebibliography}\endgroup
